\section{Markov 性质、转移矩阵与多步演化}
\label{sec:06-Random-Processes/03-Markov-Chains/01}

产品同事拿来一份会话日志：每个用户在 App 里依次访问了哪些页面——首页、搜索、商品详情、购物车、下单。他想要一个模型，给定某人已经走过的路径，预测下一步会落在哪里。

老实按定义写，要估的是

\[
P(X_{n+1}=j\given X_n,X_{n-1},\ldots,X_0).
\]

条件里挂着整段历史。页面只有十个、会话长度二十步，可能的历史就有 \(10^{20}\) 种，每一种都得配一张概率表。数据再多也填不满。

于是有人提议：只看当前页面，前面的全扔掉。这听上去像图省事的近似，其实是一句可以被数据推翻的断言——它声称当前页面已经吸干了历史里所有对预测有用的信息。本节先把这句断言写清楚，再说明它在什么时候是真的。

\subsection{Markov 性质是一句关于信息的断言}

设 \(X_n\) 是系统在离散时刻 \(n\) 的状态。Markov 性质要求

\[
P(X_{n+1}=j\given X_n=i,X_{n-1},\ldots,X_0)
=P(X_{n+1}=j\given X_n=i).
\]

若右侧还不随 \(n\) 变化，称为时间齐次 Markov 链；本节和后面几节都默认齐次。

换一种说法，这个等式说的是：\emph{给定当前状态，未来与过去条件独立}。这正是\hyperref[sec:05-Probability/02-Conditional-Probability/04]{``条件独立与信息结构''}里的链结构，过去 \(\to\) 现在 \(\to\) 未来。那一节的结论在这里原样成立：观测链条中间的节点，就掐断了两端之间的信息通道。历史仍然与未来相关，但相关性全部绕道当前状态传递；把当前状态钉住，通道就关了。

关掉这条通道换来的是压缩。原本一张以整段历史为索引的表，现在只剩十乘十的一百个数。整条路径的分布也随之写成局部规则的连乘。代价只有一个：那句条件独立的断言必须真的成立。

\subsection{状态定义得好不好，决定断言成不成立}

回到会话日志。同一个商品详情页，用户是从搜索结果点进来的，还是从购物车退回来的，下一步的去向分布并不一样——前者更可能继续浏览同类商品，后者更可能去结算。知道当前页面是``商品详情''之后，``上一页是什么''依然携带预测信息。断言被推翻了：以页面 id 为状态的链不是 Markov 链。

补救办法不是换模型，是换状态。把状态定义成``（上一页，当前页）''这个二元组，刚才泄漏出来的那点信息就被收进了状态里，链重新变成 Markov 的。参数从一百涨到一千，但仍然是常数级，而不是随会话长度指数增长。

这个手法完全一般：若下一步只依赖最近 \(k\) 步，把这 \(k\) 步打包成一个状态，\(k\) 阶依赖就化归为一阶。同一个系统，状态取得太粗就不是 Markov 的，把必要的历史并进状态就又是了。所以``这个过程是不是 Markov 的''不是关于世界的问题，是关于你怎么给状态贴标签的问题。

同样的对话在别处反复上演。一辆车的位置序列通常不 Markov，因为过去位置泄漏了速度方向；把状态改成``（位置，速度）''就补上了。排队系统若服务时间不是指数分布，只记顾客数不够，还要记当前这位已经被服务了多久，否则剩余服务时间的分布依赖于你不知道的历史。

\begin{remark}
所谓``无后效性''从来不是说世界没有历史，而是说你选的状态已经把历史压缩干净了。扩状态总能恢复 Markov 性，但状态空间会随打包的历史长度指数膨胀——这个交换是有价格的，不是免费的定理。
\end{remark}

强化学习里``MDP 的状态是否充分''问的就是这件事。智能体只看到一帧画面时，球的速度不在状态里，最优策略根本没法写成当前状态的函数；把连续几帧叠起来当状态，就是在手动补 Markov 性。序列模型的上下文窗口是同一枚硬币的另一面：不做压缩，直接把最近 \(L\) 步整段塞进``状态''。窗口一旦截断，被截掉的历史若仍有预测力，Markov 假设就在窗口边界那里破掉——长文档里指代解析出错，常常就是这么来的。

\subsection{转移矩阵把一步规则写成一张表}

有限状态空间 \(S=\set{1,\ldots,m}\) 上定义

\[
P_{ij}=P(X_{n+1}=j\given X_n=i).
\]

矩阵 \(P\) 元素非负、每一行和为一，称为行随机矩阵。行和为一是因为从 \(i\) 出发下一步总要落到某处。

\begin{center}
\begin{tikzpicture}[>=stealth,
  state/.style={draw,circle,minimum size=1.15cm,fill=blue!6,align=center},
  every node/.style={font=\small}]
  \node[state] (sun) at (-2.7,0) {晴};
  \node[state] (cloud) at (0,1.7) {阴};
  \node[state] (rain) at (2.7,0) {雨};
  \draw[->,thick] (sun) to[bend left=15] node[above left] {\(0.3\)} (cloud);
  \draw[->,thick] (cloud) to[bend left=15] node[below left] {\(0.4\)} (sun);
  \draw[->,thick] (cloud) to[bend left=15] node[above right] {\(0.2\)} (rain);
  \draw[->,thick] (rain) to[bend left=15] node[below right] {\(0.5\)} (cloud);
  \draw[->] (sun) edge[loop left] node[left] {\(0.7\)} (sun);
  \draw[->] (cloud) edge[loop above] node[above] {\(0.4\)} (cloud);
  \draw[->] (rain) edge[loop right] node[right] {\(0.5\)} (rain);
\end{tikzpicture}

\emph{每个节点的出边（含自环）概率之和为 \(1\)，这就是矩阵的行和为一；没画出的箭头对应矩阵里的零元。多步转移就是在这张图上枚举路径：沿一条路径把概率相乘，再把所有路径相加。}
\end{center}

本节把分布写成行向量。初始分布

\[
\mu_0=(P(X_0=1),\ldots,P(X_0=m))
\]

一步之后变成 \(\mu_1=\mu_0P\)。有些教材用列向量并写 \(P^{\trans}\mu\)，两种约定都自洽，混用才出事。每次乘完检查一下维度和总和是否还是一，是最省事的防线。

\subsection{多步转移就是矩阵幂}

从 \(i\) 出发两步到 \(j\)，按第一步落在哪里分情况：

\[
\begin{aligned}
P(X_{n+2}=j\given X_n=i)
&=\sum_k P(X_{n+1}=k\given X_n=i)\,P(X_{n+2}=j\given X_{n+1}=k,X_n=i)\\
&=\sum_k P_{ik}P_{kj}
=(P^2)_{ij}.
\end{aligned}
\]

第一个等号是全概率公式，第二个等号用掉了 Markov 性——正是它允许把条件里的 \(X_n=i\) 划掉；齐次性保证结果与 \(n\) 无关。归纳下去，\(n\) 步转移概率是 \((P^n)_{ij}\)，任意初始分布满足 \(\mu_n=\mu_0P^n\)。矩阵乘法在这里不是形式上的巧合，它每一次都在对所有可能的中间状态做一次全概率展开。

把路径在第 \(m\) 步而不是第一步切开，同样的论证给出 Chapman--Kolmogorov 方程 \(P^{m+n}=P^mP^n\)：先跑 \(m\) 步再跑 \(n\) 步，与一次跑 \(m+n\) 步同分布。

回到上面那张天气图。晴到雨没有直达边，矩阵里那个位置是零；但两步有通路，晴 \(\to\) 阴 \(\to\) 雨，概率 \(0.3\times0.2=0.06\)。一个元素从零变成正数，意味着图上出现了一条更长的通路——矩阵幂做的就是把图上的路径按长度分类清点。

同理，给定 \(\mu_0\)，一条具体路径的概率是

\[
P(X_0=x_0,\ldots,X_n=x_n)
=\mu_0(x_0)\prod_{k=0}^{n-1}P_{x_kx_{k+1}}.
\]

初始分布与转移矩阵合起来才决定全部有限维分布。只给 \(P\) 不给初值，过程还没被指定完。

本单元后面几节会反复用一条更小的链，状态只有晴与雨：

\[
P=\begin{bmatrix}0.8&0.2\\0.4&0.6\end{bmatrix}.
\]

它小到可以口算，却已经同时具备不可约、非周期和可逆这三样性质，正好用来把后面的抽象条件一条条对照着看。

\subsection{矩阵幂算得出答案，却解释不了答案}

原则上 \(P^n\) 包含了一切多步概率，你只要不停做乘法。但堆出来的数字不会告诉你哪些状态永远到不了、链会不会分裂成几块互不往来的区域、\(P^n\) 到底收不收敛、收敛又有多快。这些问题的答案先写在转移图的结构里，其次才写在数值里。

下一节因此暂时放下矩阵乘法，回到图上：谁能走到谁，哪些状态互相通信，以及为什么有的链会永远踩着节拍振荡而不肯安定下来。
